% Section 1: Introduction to Inverse Problems & Virtual Staining

\section{Inverse Problems in Bioimaging}

\begin{frame}[fragile]{Optical Physics: Microscope PSF Formation}
    \begin{columns}
        \column{0.48\textwidth}
        \textbf{Physics of Impulse Response}

        \begin{itemize}
            \item \textbf{Point Source}: A fluorophore emits spherical wavefronts $\delta(x,y)$.
            \item \textbf{Finite Aperture}: Objective lens with numerical aperture $\text{NA} = n \sin \alpha$ caps high spatial frequencies (low-pass filter).
            \item \textbf{Point Spread Function (PSF)}: Diffraction transforms a point source into an Airy pattern $h(r) \propto \left(\frac{J_1(k r)}{k r}\right)^2$.
            \item \textbf{Abbe Limit}: Minimum resolvable distance $d = \frac{\lambda}{2\text{NA}}$.
        \end{itemize}

        \vspace{1mm}
        \textbf{\color{commentgreen}Biological Intuition (DAPI $\to$ GFP):}\\
        \fontsize{6.5pt}{7.5pt}\selectfont
        Single GFP fluorophores or fine DAPI chromatin granules are physically blurred by the microscope's PSF into smooth Airy disks. Generative models act as non-linear deconvolution solvers.

        \column{0.52\textwidth}
        \textbf{Microscope Optics \& Code}

        \centering
        \begin{tikzpicture}[scale=0.55, every node/.style={transform shape}]
            % Axis
            \draw[dashed, gray] (-0.3,0) -- (4.8,0);
            % Object Plane
            \draw[thick, gray] (0,-1.2) -- (0,1.2) node[above, black] {\tiny Object Plane};
            \fill[red] (0,0.3) circle (2.5pt) node[left, black] {\tiny Point Source $\delta$};
            % Lens
            \draw[thick, mainblue, fill=accentblue!20] (2.0,-1.2) arc (-30:30:2.4) arc (150:210:2.4);
            \node[above, mainblue] at (2.0,1.2) {\tiny Lens (NA)};
            % Rays
            \draw[red, opacity=0.6] (0,0.3) -- (2.0,0.9) -- (4.2,-0.3);
            \draw[red, opacity=0.6] (0,0.3) -- (2.0,-0.9) -- (4.2,-0.3);
            % Image Plane & PSF Profile
            \draw[thick, gray] (4.2,-1.2) -- (4.2,1.2) node[above, black] {\tiny Image Plane};
            \fill[red] (4.2,-0.3) circle (2.5pt);
            \draw[smooth, domain=-0.9:0.9, variable=\y, color=accentblue, thick]
            plot ({4.3 + 0.6*exp(-10*(\y)^2)}, {\y - 0.3});
            \node[right, accentblue] at (4.9,-0.3) {\tiny Airy / PSF $h(r)$};
        \end{tikzpicture}

        \begin{lstlisting}
import torch

def generate_2d_gaussian_psf(size=15, sigma=2.0):
    # Coordinate grid
    coords = torch.arange(size) - size // 2
    y, x = torch.meshgrid(coords, coords, indexing='ij')
    
    # 2D Gaussian PSF approximation h(r)
    r_sq = x**2 + y**2
    psf = torch.exp(-r_sq / (2 * sigma**2))
    psf = psf / psf.sum() # Normalize integral to 1
    return psf.unsqueeze(0).unsqueeze(0)
\end{lstlisting}
    \end{columns}
\end{frame}

\begin{frame}[fragile]{Forward Problems in Bioimage Acquisition}
    \begin{columns}
        \column{0.48\textwidth}
        \textbf{Mathematical Formulation}

        Microscope image formation is a \textbf{forward problem}:
        \begin{equation*}
            y = A(x) + n
        \end{equation*}

        \begin{itemize}
            \item $x \in \R^N$: True biological signal (GFP reporter).
            \item $A$: Optical degradation (PSF blur, phase).
            \item $n \sim \mathcal{N}(0, \sigma^2 I)$: Detector noise.
            \item $y \in \R^M$: Observed micrograph (DAPI stain).
        \end{itemize}

        \vspace{1mm}
        \textbf{\color{commentgreen}Biological Intuition (DAPI $\to$ GFP):}\\
        \fontsize{6.5pt}{7.5pt}\selectfont
        DAPI ($y$) reveals nuclear boundaries and chromatin texture, which are optical projections of cell state. GFP ($x$) marks positive expression. Microscope optics blur and add noise to these signals.

        \column{0.52\textwidth}
        \textbf{Python Code Equivalent}
        \begin{lstlisting}
import torch
import torch.nn.functional as F

def forward_observation(x, psf, sigma=0.05):
    # Blur image with Point Spread Function
    blurred = F.conv2d(x, psf, padding=1)
    
    # Add Gaussian detector noise n ~ N(0, sigma^2)
    noise = torch.randn_like(blurred) * sigma
    
    y = blurred + noise
    return y
\end{lstlisting}
    \end{columns}
\end{frame}


\begin{frame}[fragile]{Ill-Posedness of the Inverse Problem}
    \begin{columns}
        \column{0.48\textwidth}
        \textbf{Mathematical Formulation}

        Reconstructing $x$ from $y$ requires solving the \textbf{inverse problem}:
        \begin{equation*}
            \hat{x} = A^{\dagger} y
        \end{equation*}

        Hadamard criteria for well-posedness:
        \begin{enumerate}
            \item \textbf{Existence}: Solution exists.
            \item \textbf{Uniqueness}: Unique solution.
            \item \textbf{Stability}: Continuous dependence on data $y$.
        \end{enumerate}

        Bioimaging inverse problems are \textbf{severely ill-posed}: $A$ is non-invertible/lossy!

        \vspace{1mm}
        \textbf{\color{commentgreen}Biological Intuition (DAPI $\to$ GFP):}\\
        \fontsize{6.5pt}{7.5pt}\selectfont
        Multiple cells may have identical nuclear shapes in DAPI, but some are GFP-positive and others negative. Inverting nuclear stain alone is ill-posed without spatial priors.

        \column{0.52\textwidth}
        \textbf{Python Code Equivalent}
        \begin{lstlisting}
# Naive inversion fails due to noise amplification

def naive_inversion(y, A_matrix):
    # SVD decomposition of measurement operator
    u, s, vh = torch.linalg.svd(A_matrix)
    
    # Small singular values amplify noise
    s_inv = torch.where(s > 1e-4, 1.0 / s, 0.0)
    A_pinv = vh.T @ torch.diag(s_inv) @ u.T
    
    x_hat = A_pinv @ y
    return x_hat
\end{lstlisting}
    \end{columns}
\end{frame}


\begin{frame}[fragile]{Classical Regularization (Tikhonov / Ridge)}
    \begin{columns}
        \column{0.48\textwidth}
        \textbf{Mathematical Formulation}

        To stabilize ill-posed inverse problems, we penalize implausible solutions:
        \begin{equation*}
            \hat{x} = \arg\min_x \|y - Ax\|_2^2 + \lambda \|x\|_2^2
        \end{equation*}

        Closed-form solution:
        \begin{equation*}
            \hat{x} = (A^T A + \lambda I)^{-1} A^T y
        \end{equation*}

        \begin{itemize}
            \item First term: Data fidelity (observed $y$).
            \item Second term: Prior/regularizer ($L_2$ norm).
            \item $\lambda > 0$: Regularization strength.
        \end{itemize}

        \vspace{1mm}
        \textbf{\color{commentgreen}Biological Intuition (DAPI $\to$ GFP):}\\
        \fontsize{6.5pt}{7.5pt}\selectfont
        Penalizes unrealistically extreme GFP intensity spikes outside cell boundaries, enforcing a smooth prior that fluorescence shouldn't explode randomly.

        \column{0.52\textwidth}
        \textbf{Python Code Equivalent}
        \begin{lstlisting}
def tikhonov_reconstruction(y, A, lmbda=1e-2):
    # A: [M, N] measurement matrix
    # y: [M, 1] observation vector
    N = A.shape[1]
    
    # System: (A^T A + lambda * I) x = A^T y
    AtA = A.T @ A
    regularizer = lmbda * torch.eye(N, device=A.device)
    
    lhs = AtA + regularizer
    rhs = A.T @ y
    
    x_hat = torch.linalg.solve(lhs, rhs)
    return x_hat
\end{lstlisting}
    \end{columns}
\end{frame}


\begin{frame}[fragile]{Autoencoders (AEs) \& Representation Learning}
    \begin{columns}
        \column{0.48\textwidth}
        \textbf{Mathematical Formulation}

        Compression into a low-dimensional bottleneck $z \in \R^d$ ($d \ll N$):
        \begin{align*}
            z &= e_\phi(x) \quad \text{(Encoder)} \\
            \hat{x} &= d_\theta(z) \quad \text{(Decoder)}
        \end{align*}
        Reconstruction Loss: $\mathcal{L}_{\text{AE}} = \|x - d_\theta(e_\phi(x))\|_2^2$

        \textbf{Generative Sampling Limitation}:
        Deterministic $z$-space has gaps and unconstrained regions. Random sampling $z \sim \mathcal{N}(0, I)$ produces unrealistic artifacts!

        \vspace{1mm}
        \textbf{\color{commentgreen}Biological Intuition (DAPI $\to$ GFP):}\\
        \fontsize{6.5pt}{7.5pt}\selectfont
        Standard AE compresses DAPI nuclear shape and GFP expression into bottleneck $z$. But because $z$-space has gaps, sampling random vectors yields non-biological cell artifacts.

        \column{0.52\textwidth}
        \textbf{Python Code Equivalent}
        \begin{lstlisting}
import torch.nn as nn

class Autoencoder(nn.Module):
    def __init__(self, latent_dim=32):
        super().__init__()
        self.encoder = nn.Sequential(
            nn.Conv2d(1, 16, 3, stride=2, padding=1),
            nn.ReLU(),
            nn.Flatten(),
            nn.Linear(16 * 128 * 128, latent_dim)
        )
        self.decoder = nn.Sequential(
            nn.Linear(latent_dim, 16 * 128 * 128),
            nn.Unflatten(1, (16, 128, 128)),
            nn.ConvTranspose2d(16, 1, 3, stride=2, padding=1, output_padding=1),
            nn.Sigmoid()
        )

    def forward(self, x):
        z = self.encoder(x)
        return self.decoder(z)
\end{lstlisting}
    \end{columns}
\end{frame}


\begin{frame}[fragile]{Variational Autoencoders \& Generative Sampling}
    \begin{columns}
        \column{0.48\textwidth}
        \textbf{Mathematical Formulation}

        Map input $x$ to latent distribution parameters $(\mu_\phi(x), \sigma_\phi(x))$:
        \begin{equation*}
            z = \mu_\phi(x) + \sigma_\phi(x) \odot \epsilon, \quad \epsilon \sim \mathcal{N}(0, I)
        \end{equation*}
        Objective (ELBO):
        \begin{equation*}
            \mathcal{L}_{\text{VAE}} = \|x - \hat{x}\|_2^2 + \beta \, \KL{q_\phi(z|x)}{\mathcal{N}(0, I)}
        \end{equation*}

        \textbf{Generative Sampling}:
        $\text{KL}$ penalty regularizes $z$-space to be smooth and gap-free. We can synthesize NEW valid cell images by sampling $z \sim \mathcal{N}(0, I)$!

        \vspace{1mm}
        \textbf{\color{commentgreen}Biological Intuition (DAPI $\to$ GFP):}\\
        \fontsize{6.5pt}{7.5pt}\selectfont
        VAE regularizes cell latent space so nuclear features and GFP expression vary smoothly. We can sample new synthetic DAPI/GFP cells or interpolate between GFP-negative and GFP-positive states!

        \column{0.52\textwidth}
        \textbf{Python Code Equivalent}
        \begin{lstlisting}
# Generative Sampling from Trained VAE Decoder

def generate_samples(decoder, num_samples=16, latent_dim=32, device='cpu'):
    decoder.eval()
    with torch.no_grad():
        # 1. Sample z ~ N(0, I) from standard Gaussian prior
        z_random = torch.randn(num_samples, latent_dim, device=device)
        
        # 2. Decode latent vectors into synthetic cell images
        x_synthetic = decoder(z_random)
        
    return x_synthetic # [num_samples, 1, H, W]
\end{lstlisting}
    \end{columns}
\end{frame}


\begin{frame}[fragile]{Learning the Prior: Deep Inverse Problems}
    \begin{columns}
        \column{0.48\textwidth}
        \textbf{Mathematical Formulation}

        Instead of hand-crafted priors $\|x\|_2^2$ or $\text{TV}(x)$, we learn a data-driven mapping $f_\theta$:
        \begin{equation*}
            \hat{x} = f_\theta(y) \approx \arg\min_x \Big( \|y - A(x)\|^2 + \lambda R_\theta(x) \Big)
        \end{equation*}

        Empirical Risk Minimization over training pairs $(y_i, x_i)$:
        \begin{equation*}
            \theta^* = \arg\min_\theta \frac{1}{N} \sum_{i=1}^N \mathcal{L}\big( f_\theta(y_i), \, x_i \big)
        \end{equation*}

        \vspace{1mm}
        \textbf{\color{commentgreen}Biological Intuition (DAPI $\to$ GFP):}\\
        \fontsize{6.5pt}{7.5pt}\selectfont
        A deep neural network learns subtle nuclear features (e.g., nuclear swelling or DAPI texture changes) to predict whether a cell is GFP-positive or negative.

        \column{0.52\textwidth}
        \textbf{Python Code Equivalent}
        \begin{lstlisting}
import torch.nn as nn

class DeepInverseSolver(nn.Module):
    def __init__(self, backbone_unet):
        super().__init__()
        self.unet = backbone_unet # f_theta
        
    def forward(self, y):
        # Directly map degraded image y to reconstructed x
        return self.unet(y)

# Training objective over paired dataset
def compute_reconstruction_loss(x_pred, x_gt):
    return F.l1_loss(x_pred, x_gt)
\end{lstlisting}
    \end{columns}
\end{frame}


\begin{frame}[fragile]{Virtual Staining as a Non-Linear Inverse Problem}
    \begin{columns}
        \column{0.48\textwidth}
        \textbf{Mathematical \& VIRVS Formulation}

        \begin{itemize}
            \item \textbf{Input ($y$)}: DAPI/Hoechst nuclear micrograph.
            \item \textbf{Target ($x$)}: Fluorescent GFP reporter ($>0$ for positive cell, $0$ for negative).
            \item \textbf{Inverse Task} (Wyrzykowska et al., 2024):
        \end{itemize}

        \begin{equation*}
            x_{\text{GFP}} \approx f_\theta( y_{\text{DAPI}} )
        \end{equation*}

        \vspace{1mm}
        \textbf{\color{commentgreen}Biological Intuition (DAPI $\to$ GFP):}\\
        \fontsize{6.5pt}{7.5pt}\selectfont
        Predicts continuous GFP reporter intensity per cell directly from nuclear DAPI staining, enabling non-invasive quantification of positive vs negative cells!

        \column{0.52\textwidth}
        \textbf{Python Code Equivalent}
        \begin{lstlisting}
# Virus Infection Reporter Virtual Staining (VIRVS)
# Mapping Brightfield/DAPI y -> GFP Reporter x

def virvs_forward_step(model, dapi_img, gfp_target, optimizer):
    optimizer.zero_grad()
    
    # Predict continuous GFP fluorescence x_hat
    gfp_pred = model(dapi_img) # f_theta(y)
    
    # Mean Absolute Error Loss
    loss = F.l1_loss(gfp_pred, gfp_target)
    loss.backward()
    optimizer.step()
    
    return loss.item(), gfp_pred
\end{lstlisting}
    \end{columns}
\end{frame}
